\documentclass[twocolumn]{aastex631}
\graphicspath{{../figures/}}
\begin{document}

\title{SDSS high-excitation AGN denominator for outflow escape tests: an SDSS DR17 pilot}
\shorttitle{SDSS high-excitation AGN denominator for outflow escape tests}
\author{NebulaMind Research Autopilot}
\affiliation{Local reproducible pilot run; public SDSS DR17 data only}

\begin{abstract}
We execute a bounded actual-data pilot for the Galaxy Evolution research topic ``Escape versus recycling: the fate of AGN-driven multiphase outflows''.  The analysis uses a public SDSS DR17 emission-line galaxy sample with stellar-mass, specific-SFR, photometry, and optical emission-line measurements.  The pilot question is: How large is the SDSS high-excitation optical-AGN denominator that would need resolved kinematics to test escape versus recycling?  The result is intended as a survey denominator or proxy measurement, not as a full causal test of feedback physics.
\end{abstract}

\section{Scope}
This manuscript is an AAS-style pilot generated from a research proposal.  It uses actual public survey data and preserves the analysis artifacts, but it deliberately stays inside the observables available in the SDSS sample.  The full proposal requires resolved outflow velocities, halo potentials, molecular/ionized/neutral gas phases, and CGM recycling tracers.

\section{Data and Measurements}
The input table is the SDSS DR17 emission-line sample from run SDSS-AGN-SFR-PILOT-20260708T122000Z.  It contains 60,000 galaxies after requiring spectroscopic galaxy classification, redshift 0.02--0.12, finite stellar-mass and specific-SFR estimates, and signal-to-noise at least 3 in H$\alpha$, H$\beta$, [O~III] $\lambda5007$, and [N~II] $\lambda6584$.  BPT classes are recomputed from the line ratios using the Kauffmann et al. and Kewley et al. demarcations.  A local-density ranking is computed from the 10th nearest neighbour in approximate comoving Cartesian coordinates and is used only as an internal density proxy.

\section{Pilot Result}
\begin{itemize}
\item High-excitation optical AGN candidates number 4,440 of 60,000 emission-line galaxies (0.074).
\item Their median log sSFR is -11.53, compared with -10.14 for the full denominator.
\item SDSS does not measure escape velocity or multiphase outflow velocities here; the pilot supplies a denominator for resolved follow-up rather than an escape/recycling result.
\end{itemize}

\begin{figure}
\centering
\includegraphics[width=\columnwidth]{m2\_p1\_outflow\_escape\_recycling\_figure1.pdf}
\caption{Topic-specific SDSS DR17 pilot measurement. The figure is a proxy or denominator diagnostic, not the full multi-survey test.}
\end{figure}

\section{Interpretation Guard}
This pilot should not be read as a completed proof of the full feedback proposal.  It is a reproducible SDSS measurement that defines a denominator, proxy, or validation target for the larger research design.  In particular, it does not by itself establish causal AGN feedback, gas escape, radio-jet coupling efficiency, X-ray maintenance heating, molecular-gas depletion, or simulation correctness.

\section{Reproducibility}
Run identifier: SDSS\_REMAINING\_TOPIC\_PILOTS\_20260708T125828Z.  The run directory preserves the topic-specific JSON summary, figure files, manuscript source, compile log, and compiled PDF.  The workflow used read-only public SDSS-derived data already cached from the first pilot plus local artifact writes only.

\acknowledgments
This pilot used public SDSS DR17 data products and open-source Python tools.

\begin{thebibliography}{}
\bibitem[Baldwin et al.(1981)]{baldwin1981} Baldwin, J.~A., Phillips, M.~M., \& Terlevich, R. 1981, PASP, 93, 5
\bibitem[Kauffmann et al.(2003)]{kauffmann2003} Kauffmann, G., Heckman, T.~M., Tremonti, C., et al. 2003, MNRAS, 346, 1055
\bibitem[Kewley et al.(2001)]{kewley2001} Kewley, L.~J., Dopita, M.~A., Sutherland, R.~S., Heisler, C.~A., \& Trevena, J. 2001, ApJ, 556, 121
\bibitem[York et al.(2000)]{york2000} York, D.~G., Adelman, J., Anderson, J.~E., Jr., et al. 2000, AJ, 120, 1579
\end{thebibliography}

\end{document}
